
\subsection{Valuations: positive and negative cases, and the remaining test}

\begin{claim}{FF.16}{Proved}\label{clm:FFmixed}
Residue characteristic \(p\) does not, even together with any
specified finite transcendence degree over \(\mathbb Q\), decide
reconstruction. For each prime \(p\) and integer \(d\geq0\),
there exist henselian valued fields of transcendence degree \(d\)
over \(\mathbb Q\), with residue characteristic \(p\), which
reconstruct, and others which fail reconstruction.
\end{claim}
\begin{proof}
Choose algebraically independent elements
\(a_1,\ldots,a_d\in\mathbb Q_p\), and give
\(D=\mathbb Q(a_1,\ldots,a_d)\) the restricted \(p\)-adic
valuation. Such a tuple exists: \(\mathbb Q_p\) is uncountable,
whereas a field of countable transcendence degree over the
countable field \(\mathbb Q\), and its algebraic closure, are
countable. Henselization universality embeds \(D^h\) as a
valued subfield of the henselian field \(\mathbb Q_p\).
Henselization is algebraic and immediate
\citep[Proposition~3.3.25, Corollary~3.3.26]{FF_ADH2017}.
Thus \(\operatorname{trdeg}_{\mathbb Q}D^h=d\), its value
group is \(\mathbb Z\), and its residue field is
\(\mathbb F_p\): these assertions hold for \(D\) because
\(\mathbb Q\subset D\subset\mathbb Q_p\).
The field \(D^h\) is sub-\(p\)-adic, so Carlson--Stix
Theorem~A gives \(\mathcal R(D^h)\)
\citep[Theorem~A, Theorem~4.1]{FF_CS2026}.
For \(d=0\) this is already the ordinary henselization of
\(\mathbb Q\) at \(p\), embedded in \(\mathbb Q_p\).

For the opposite verdict take an algebraic closure
\(K=\overline{\mathbb Q(t_1,\ldots,t_d)}\).
Extend the \(p\)-adic Gauss valuation on the displayed rational
function field to \(K\). Its residue characteristic remains
\(p\), and an algebraically closed valued field is henselian.
The transcendence degree is \(d\). The group \(G_K\) is
trivial while \(\operatorname{Aut}(K)\) is nontrivial; a
nonidentity automorphism of \(\overline{\mathbb Q}\) extends
after choosing a transcendence basis. The centralizer obstruction
of Claim~\ref{clm:FF6} therefore gives \(\neg\mathcal R(K)\).
This supplies both verdicts without inferring a uniform answer
from the residue characteristic.
\end{proof}

\begin{claim}{FF.17}{Proved}\label{clm:FFdefect}
There is an explicit henselian field \(K/\mathbb Q\) of
transcendence degree two, residue characteristic zero and
transcendence defect one which fails reconstruction. Namely, in
\(\overline{\mathbb Q}((t))\) put
\[
 u=\sum_{n\geq1}t^{n!},\qquad
 K=\overline{\mathbb Q}(t,u)^h
\]
for the induced \(t\)-adic valuation and its henselization
inside that complete field. Then
\[
 \kappa(K)=\overline{\mathbb Q},\qquad
 v(K^\times)=\mathbb Z,\qquad G_K\simeq\widehat{\mathbb Z}.
\tag{FF.17.1}
\]
The transcendence degree two is the smallest possible finite
transcendence degree for a nontrivial residue-characteristic-zero
valuation with positive transcendence defect.
\end{claim}
\begin{proof}
The power series \(u\) is transcendental over
\(\overline{\mathbb Q}(t)\). Indeed let
\(P(t,U)\in\overline{\mathbb Q}[t,U]\) be nonzero, of
\(U\)-degree \(e\) and \(t\)-degree at most \(m\), and
write \(u_N=\sum_{n=1}^N t^{n!}\). If \(P(t,u)=0\), then
\[
 v\bigl(P(t,u_N)\bigr)\geq(N+1)!,\qquad
 \deg_tP(t,u_N)\leq m+eN!
\]
whenever the polynomial on the left is nonzero. For all
sufficiently large \(N\) these inequalities are incompatible.
Consequently infinitely many distinct \(u_N\) would be roots
of the nonzero polynomial \(P(t,U)\) over
\(\overline{\mathbb Q}(t)\), a contradiction.
Henselization is algebraic and immediate, proving the asserted
transcendence degree, residue field and value group.

For compatible roots of \(t\), let
\(E=\bigcup_n K(t^{1/n})\). Each extension has degree \(n\)
by the distinct valuation classes of
\(1,t^{1/n},\ldots,t^{(n-1)/n}\); since all roots of unity
are in \(K\), it is cyclic Galois. The union is henselian,
with residue field \(\overline{\mathbb Q}\) and divisible
value group \(\mathbb Q\). In residue characteristic zero,
these two conditions make a henselian field algebraically closed
\citep[Corollary~3.3.22]{FF_ADH2017}. Thus \(E=\bar K\)
and \(G_K=\varprojlim_n\mu_n\simeq\widehat{\mathbb Z}\).
Its nontrivial centre contradicts the necessary condition in
Claim~\ref{clm:FF6}. Notice that this obstruction may fix \(K\)
pointwise: it is a nontrivial central automorphism of its algebraic
closure, detected by the retained relative structural 2-cell.

Finally the rational rank of a nontrivial ordered value group
is at least one. The residue transcendence degree is nonnegative.
Positive transcendence defect therefore requires
\(\operatorname{trdeg}_{\mathbb Q}K\geq2\). In our example
\(2-0-1=1\), as claimed. The attached finite script checks
the displayed factorial-gap inequalities for sample bounds; the
argument above proves them for all bounds and proves transcendence.
\end{proof}

\begin{claim}{FF.18}{Open}\label{clm:FFremaining}
The following sufficient direction of candidate F3(a) remains
unproved and unrefuted in this work:
\[
 \bigl[C_{\operatorname{Aut}(\bar k)}(G_E)=1
          \text{ for every finite }E/k\bigr]
                  \quad\Longrightarrow\quad\mathcal R(k).
\tag{FF.18.1}
\]
Nor have we decided every finite-transcendence-degree henselian
field with residue characteristic zero and transcendence defect.
In particular our arguments do not decide
\[
 K_* = \mathbb Q\bigl(t,\sum_{n\geq1}t^{n!}\bigr)^h
           \subset\mathbb Q((t)),
\tag{FF.18.2}
\]
with its induced valuation. This is a specified unresolved test,
not a claim that it is the least field under inclusion or that
its centralizers satisfy (FF.18.1).
\end{claim}
\begin{proof}[State of the attempted proof and its exact boundary]
Claim~\ref{clm:FF6} proves the necessary direction. The known
substitution ghosts of AP-F.41, F.42, F.47 and F.48 fail its
hypothesis, so none refutes the converse. In particular one must
not infer sufficiency from having checked all the previous
counterexamples. Attempts to replace the reconstruction argument
by this rigidity condition have not supplied the missing step:
centralizers control invertible semilinear ghosts, whereas a
general pinned map to \(\mathbb G_m\) has Kummer cocycles and
actual point values which have not been forced to descend by
centralizer triviality alone. The finite-quotient detection and
normal arithmetic action in Claim~\ref{clm:FFcore} are genuine
extra inputs; we have not deduced them from (FF.18.1).

For valuations, AP-F.48 supplies a coefficient field \(F\) and
parameters with rationally independent values. Under the
Abhyankar equality the ambient field is algebraic over the
corresponding monomial henselization, so algebraic descent
transfers its failure. Positive transcendence defect removes
that algebraicity statement. It does not remove all failures,
as Claim~\ref{clm:FFdefect} shows, but that example uses
algebraically closed residue field and a central Galois group.
For (FF.18.2) the residue field is \(\mathbb Q\), and the
same radical-and-constant argument gives
\(G_{K_*}\simeq\widehat{\mathbb Z}(1)\rtimes G_{\mathbb Q}\).
Its centre does not supply the previous obstruction.

There is a concrete obstruction to copying the old substitution
proof. Put \(u=\sum t^{n!}\). Then
\(v(u-t-t^2)=6\), whereas substitution
\(t\mapsto t+t^2\), \(u\mapsto u\), gives
\[
 u-(t+t^2)-(t+t^2)^2
       =-t^2-2t^3-t^4+t^6+\cdots,
\]
of valuation two. Thus that field embedding of
\(\mathbb Q(t,u)\) is not valuation preserving and does not
extend by the henselization universal property used in AP-F.48.
Allowing \(u\) to move would require an actual algebraic
self-correspondence preserving its formal branch; none has been
constructed here. Attempts to use rigid henselian fields would
also require proving the centralizer condition for every finite
extension and independently constructing a nonalgebraic pinned
map. Neither assertion follows from field rigidity alone.

Similarly, Claim~\ref{clm:FFpro} proves the pro-form for the
explicit coherent relative model category. The transition maps
of the models need not be pinned, so simply writing an ordinary
\(\operatorname{Pro}(\operatorname{Top}^{\mathrm{pin}})\)
is invalid. We have not proved that any other proposed definition
of a closed-point-preserving ordinary pro-morphism has the same
morphisms as that coherent category. The literal unspecified
pro-category reading of F1(iii), and that part of F6, is therefore
not silently marked proved. The correct coherent-descent version
of F3(c) is Claim~\ref{clm:FFc}; the bare commuting-square
version is refuted by Claim~\ref{clm:FF8}.
\end{proof}
